\section{A uniform estimate for cell moments}
\label{sec:moments}

We compare changes in the first moments and masses of cells when their
intercepts vary and their sites remain fixed. The estimate is uniform
over all configurations with positive cell masses.

\begin{lemma}[Stationary-site moment estimate]\label{moments:derivative}
Fix distinct sites $y_1,\ldots,y_N$ and intercepts $c_1,\ldots,c_N$
whose cells have positive masses $M_i$. If $K\subset\{x:|x|\le R\}$,
then every intercept direction $q\in\R^N$ satisfies
\[
  \sum_{i=1}^N\frac{|D B_i[q,0]|^2}{M_i}
  \le C_0R^2\sum_{i=1}^N\frac{|D M_i[q,0]|^2}{M_i},
  \qquad C_0=162.
\]
\end{lemma}

The proof couples uniform laws on two opposite perturbations of each
cell. Their midpoint lies in the original cell; the following lemma
controls its density.

\begin{lemma}[Midpoint coupling]\label{moments:midpoint}
Let $U,V\subset\R^d$ be compact convex sets of positive volume. There is a
coupling $(X,Z)$ of the uniform probability measures on $U$ and $V$ such that
\[
  \operatorname{Law}\!\left(\frac{X+Z}{2}\right)
  \le \frac{dx}{\sqrt{|U||V|}}.
\]
The inequality is an inequality of measures on $\R^d$.
\end{lemma}

This density bound also follows from McCann's displacement-concavity
inequality \cite[Theorem~2.3]{McCann1997}. We give a direct proof by the
conditional-quantile construction underlying Knothe's triangular
rearrangement \cite{Knothe1957,CarlierGalichonSantambrogio2010}.
The subsequent application to cell moments is developed here.

\begin{proof}
We parameterize the two uniform laws by successive one-dimensional
quantiles on the same cube. Their diagonal derivatives have constant
products. Averaging the two maps and substituting coordinates successively
will give the density bound.

We need the boundary observation from the proof of \Cref{cells:calculus}
for arbitrary positive-dimensional affine slices, not only hyperplanes.
The same segment-extension argument shows that a relative interior point
of a slice of $U$ meeting $\operatorname{int}U$ is an ambient interior
point: unless it equals a chosen ambient interior point in the slice,
it lies strictly between that point and a farther point of the section.
The contraction argument for convex
boundaries, applied within the slice, then shows that its intersection
with $\partial U$ has zero slice volume.

For $0\le k<d$, let
\[
  p_k(\xi)=\int_{\R^{d-k}}
       1_{\operatorname{int}U}(\xi,\eta)\,d\eta,
  \qquad
  p_d=1_{\operatorname{int}U}.
\]
Thus $p_0=|U|$. Write $D_k$ for the projection of
$\operatorname{int}U$ onto its first $k$ coordinates. Each $p_k$ is
positive and continuous on $D_k$. For $k<d$, continuity follows by
dominated convergence: at a prefix in $D_k$, the preceding boundary
property makes the indicators of the fibers converge almost everywhere
as the prefix varies, and boundedness of $U$ provides a common integrable
bound. For $k=d$ the assertion is immediate. The fiber integrals are
Borel functions on their full domains.

Given $s=(s_1,\ldots,s_d)\in(0,1)^d$, construct $X(s)$ recursively.
For a previously constructed prefix
$\xi=(X_1,\ldots,X_{k-1})\in D_{k-1}$, let $X_k$ be the quantile
at level $s_k$ of the density
\[
  t\longmapsto \frac{p_k(\xi,t)}{p_{k-1}(\xi)}.
\]
Fubini's theorem shows that this density integrates to one. It is
positive and continuous on the bounded open interval
$\{t:(\xi,t)\in D_k\}$, and vanishes outside that interval.
Its cumulative distribution function is therefore strictly increasing
from zero to one there. This defines the quantile uniquely, keeps the
new prefix in $D_k$, and gives
\[
  \partial_{s_k}X_k
   =\frac{p_{k-1}(X_1,\ldots,X_{k-1})}
          {p_k(X_1,\ldots,X_k)}>0.
\]
For each fixed earlier parameter prefix, $X_k$ is $C^1$ in $s_k$.
The conditional cumulative distribution functions are Borel in their
arguments, so comparison at rational points shows that their quantiles
are Borel. Thus the whole construction is measurable. With $s$ uniform
on the cube, successive conditioning gives the uniform law on $U$.
The diagonal derivatives telescope to
\[
  \prod_{k=1}^d\partial_{s_k}X_k=|U|.
\]
Construct $Z(s)$ in the same way for $V$, using the same parameters.
For $H=(X+Z)/2$, the arithmetic--geometric mean inequality yields
\[
  \prod_{k=1}^d\partial_{s_k}H_k
  \ge
  \sqrt{\left(\prod_{k=1}^d\partial_{s_k}X_k\right)
              \left(\prod_{k=1}^d\partial_{s_k}Z_k\right)}
  =\sqrt{|U||V|}.
\]

We justify substitution using only the stated regularity in each
coordinate. For a nonnegative Borel function $g$ on $\R^d$, fix
$a=(s_1,\ldots,s_{d-1})$. One-dimensional substitution in the strictly
increasing $C^1$ function $s_d\mapsto H_d(a,s_d)$, followed by enlargement
of its image interval to $\R$, gives
\[
  \int_0^1 g\bigl(H_1(a),\ldots,H_{d-1}(a),H_d(a,t)\bigr)
                 \partial_{s_d}H_d(a,t)\,dt
  \le \int_{\R}g\bigl(H_1(a),\ldots,H_{d-1}(a),z\bigr)\,dz.
\]
Here each $H_k(a)$ uses only the first $k$ entries of $a$.
The integral of $g$ in its last variable is again a nonnegative Borel
function. The earlier derivative factors depend only on earlier
parameters. Tonelli's theorem therefore permits multiplication by their
product and integration in $a$. Repeating this argument backwards gives
\[
  \int_{(0,1)^d}g(H(s))
       \prod_{k=1}^d\partial_{s_k}H_k(s)\,ds
  \le \int_{\R^d}g(x)\,dx.
\]
All derivative factors are Borel, as also follows from their explicit
formulas above. No differentiability in earlier parameters or boundedness
of the derivatives near the boundary of the cube is required.
Combining the last two inequalities proves
\[
  \E g\!\left(\frac{X+Z}{2}\right)
  \le \frac{1}{\sqrt{|U||V|}}\int_{\R^d}g(x)\,dx,
\]
which is the asserted measure domination.
\end{proof}

\begin{proof}[Proof of \Cref{moments:derivative}]
Fix the intercept direction $q$ and hold the sites fixed.
For sufficiently small $s>0$, let $P_i^-$ and
$P_i^+$ be the cells at $c-sq$ and $c+sq$, respectively, and write
\[
  \alpha_i=\rho(P_i^-),\qquad
  \beta_i=\rho(P_i^+),\qquad
  \gamma_i=\sqrt{\alpha_i\beta_i}.
\]
These masses are positive by continuity from \Cref{cells:calculus}.
The perturbed cells are compact convex sets of positive volume.
Couple their uniform probability measures by \Cref{moments:midpoint},
writing $X_i\in P_i^-$, $Z_i\in P_i^+$, and
$D_i=\E|X_i-Z_i|^2$. For each competing index $j$, the defining
inequalities give
\begin{align*}
  X_i\cdot(y_i-y_j)+c_i-c_j-s(q_i-q_j)&\ge0,\\
  Z_i\cdot(y_i-y_j)+c_i-c_j+s(q_i-q_j)&\ge0.
\end{align*}
Averaging shows that $(X_i+Z_i)/2$ satisfies the $i$-cell inequalities
at $c$; convexity also places it in $K$. Its law is supported on $P_i$.
Since $|P_i^-|=|K|\alpha_i$ and $|P_i^+|=|K|\beta_i$, the midpoint
bound has exactly the normalization
\[
  \gamma_i\operatorname{Law}\!\left(\frac{X_i+Z_i}{2}\right)
  \le \frac{1_{P_i}\,dx}{|K|}=\rho|_{P_i}.
\]

Use the nonnegative function $\Phi(x)=R^2-|x|^2$ on $K$, and set
$A_i^-=\E\Phi(X_i)$ and $A_i^+=\E\Phi(Z_i)$. The identity
\[
  \Phi\!\left(\frac{X_i+Z_i}{2}\right)
  =\frac{\Phi(X_i)+\Phi(Z_i)}2+\frac{|X_i-Z_i|^2}{4}
\]
and the cellwise measure domination imply
\[
  \sum_i\gamma_i\left(\frac{A_i^-+A_i^+}{2}+\frac{D_i}{4}\right)
  \le\int_K\Phi\,d\rho
  =\frac12\sum_i(\alpha_iA_i^-+\beta_iA_i^+).
\]
The last equality uses the partitions at the two perturbed intercepts;
all three partitions have null overlaps because the sites are distinct.

Put $\delta_i=\alpha_i-\beta_i$ and
\[
  \mathcal A=\sum_i\gamma_iD_i,
  \qquad
  \mathcal B=\sum_i\frac{\delta_i^2}{\gamma_i}.
\]
Subtracting the averaged endpoint terms in the preceding inequality
gives
\[
  \frac{\mathcal A}{4}
  \le\sum_i\left[
    \left(\frac{\alpha_i+\beta_i}{2}-\gamma_i\right)
        \frac{A_i^-+A_i^+}{2}
    +\frac{\delta_i}{4}(A_i^--A_i^+)
  \right].
\]
Here $0\le A_i^\pm\le R^2$, and
\[
  \frac{\alpha_i+\beta_i}{2}-\gamma_i
  =\frac{\delta_i^2}{2(\sqrt{\alpha_i}+\sqrt{\beta_i})^2}
  \le\frac{\delta_i^2}{8\gamma_i}.
\]
Moreover, since $|X_i|,|Z_i|\le R$,
\[
  |A_i^--A_i^+|
  \le 2R\E|X_i-Z_i|\le2R\sqrt{D_i}.
\]
Cauchy--Schwarz therefore yields
\[
  \frac{\mathcal A}{4}
  \le\frac{R^2\mathcal B}{8}
       +\frac R2\sqrt{\mathcal A\mathcal B}.
\]
Using
$\frac R2\sqrt{\mathcal A\mathcal B}
 \le\mathcal A/8+R^2\mathcal B/2$
and rearranging, we obtain
\[
  \mathcal A\le5R^2\mathcal B.
\]
This argument includes the cases where either quantity vanishes.

The cell-moment difference is
\[
  \Delta B_i:=B_i(c-sq)-B_i(c+sq)
       =\alpha_i\E X_i-\beta_i\E Z_i,
\]
where the notation suppresses the fixed sites. Consequently,
\[
  |\Delta B_i|
  \le R|\delta_i|+\alpha_i\sqrt{D_i}.
\]
For this fixed finite configuration and fixed direction $q$, continuity
allows $s$ to be chosen small enough that every
$\alpha_i,\beta_i,\gamma_i$ lies between $M_i/2$ and $2M_i$.
Set $\mathcal E=\sum_i\delta_i^2/M_i$. These comparisons give
$\alpha_i^2/M_i\le8\gamma_i$ and $\mathcal B\le2\mathcal E$.
It follows that
\begin{align*}
  \sum_i\frac{|\Delta B_i|^2}{M_i}
  &\le2R^2\mathcal E+2\sum_i\frac{\alpha_i^2}{M_i}D_i\\
  &\le2R^2\mathcal E+16\mathcal A\\
  &\le162R^2\mathcal E.
\end{align*}
Divide by $(2s)^2$ and let $s$ decrease to zero.
By \Cref{cells:calculus},
$\Delta B_i/(2s)\to-D B_i[q,0]$ and
$\delta_i/(2s)\to-D M_i[q,0]$, which proves the claim.

The permitted threshold for $s$ may depend on the direction and the
finite configuration, including its smallest mass. Only the uniform
constant $162$ remains after taking the derivative limit. Thus the
estimate has no dependence on the number of cells, their masses, or
the separation of the sites.
\end{proof}
